\begin{figure}[htbp]
  \centering
  \begin{tikzpicture}[fig, node distance=6mm,
      big/.style={proc, text width=48mm}]
    \node[proc, text width=30mm, fill=white] (u) at (0,0) {\textbf{Learner} $u$};
    \node[big] (inf) at (5.4,1.55) {\textbf{Infer}\\[1pt]{\scriptsize update $P(L_{u,s})$ for the answered
      skill, then propagate along the prerequisite graph}};
    \node[big] (sel) at (5.4,-1.55) {\textbf{Select}\\[1pt]{\scriptsize choose the next skill at the
      knowledge frontier, a Bloom level one band above, and a question}};
    \node[store, text width=30mm] (graph) at (11.3,1.55) {skill graph $(S,E)$\\{\scriptsize acyclic, depth 9}};
    \node[store, text width=30mm] (bank) at (11.3,-1.55) {question bank\\{\scriptsize tagged with skill, Bloom, topic}};
    \draw[flowHi] (u.north) |- node[lbl, pos=0.75, above]{response $(q_t, c_t)$} (inf.west);
    \draw[flowHi] (inf.south) -- node[lbl, right=1pt]{posterior over all $s \in S$} (sel.north);
    \draw[flowHi] (sel.west) -| node[lbl, pos=0.25, below]{next question $q_{t+1}$} (u.south);
    \draw[flow] (graph) -- (inf);
    \draw[flow] (bank) -- (sel);
    \node[note, text width=40mm, anchor=north] at (5.4,-3.0)
      {output to the learner: a mastery map --- which skill, specifically, is missing};
  \end{tikzpicture}
  \caption[The problem as a closed inference--selection loop]{The problem as a closed loop. Each response is weak, noisy evidence
    about one skill; the inference step turns it into a posterior over every skill
    in the graph, and the selection step spends the next question where that
    posterior is least certain. The two stores are the fixed inputs defined in
    \Cref{sec:problem-statement}.}
  \label{fig:adaptive-loop}
\end{figure}
